\documentclass[aspectratio=169,]{beamer}

% -------------------------------------------------------------------------
% 1. REQUIRED PACKAGES (The "Batteries" Quarto usually includes)
% -------------------------------------------------------------------------
\usepackage{graphicx}
\usepackage{booktabs}
\usepackage{longtable}
\usepackage{hyperref}
\usepackage{array}      % <--- ADD THIS LINE HERE
\usepackage[most]{tcolorbox} % REQUIRED for Quarto Callouts
\usepackage{fontawesome5}    % REQUIRED for Callout Icons
\usepackage{framed}          % REQUIRED for standard blocks
\usepackage{tikz}
\usetikzlibrary{arrows.meta}

% Fix for Pandoc's "tightlist"
\providecommand{\tightlist}{%
  \setlength{\itemsep}{0pt}\setlength{\parskip}{0pt}}

% -------------------------------------------------------------------------
% 2. CUSTOM DESIGN PACKAGES
% -------------------------------------------------------------------------
\usepackage{tikz}
\usepackage{calc}
\usepackage{xcolor}
\usepackage{xparse} 
\usepackage{etoolbox} 
\usepackage[default]{sourcesanspro} 
\usetikzlibrary{calc, positioning, shapes.geometric}

% -------------------------------------------------------------------------
% DEFINE QUARTO CALLOUT COLORS (REQUIRED)
% -------------------------------------------------------------------------
% We map these to your "Elegant Dark Mode" aesthetic.
% Note: Quarto expects these EXACT names.

% 1. Note (Blue)
\definecolor{quarto-callout-note-color}{HTML}{3498DB}        % Icon/Title Color
\definecolor{quarto-callout-note-color-frame}{HTML}{2980B9}  % Border Color

% 2. Tip (Green)
\definecolor{quarto-callout-tip-color}{HTML}{27AE60}
\definecolor{quarto-callout-tip-color-frame}{HTML}{2ECC71}

% 3. Important (Red/Crimson)
\definecolor{quarto-callout-important-color}{HTML}{C0392B}
\definecolor{quarto-callout-important-color-frame}{HTML}{E74C3C}

% 4. Warning (Orange)
\definecolor{quarto-callout-warning-color}{HTML}{F39C12}
\definecolor{quarto-callout-warning-color-frame}{HTML}{E67E22}

% 5. Caution (Yellow/Gold)
\definecolor{quarto-callout-caution-color}{HTML}{F1C40F}
\definecolor{quarto-callout-caution-color-frame}{HTML}{D4AC0D}

% Add this right before \begin{document}
\newcommand{\Est}[1]{\textcolor{quarto-callout-important-color}{\mathbf{#1}}}


% -------------------------------------------------------------------------
% NEW: COLOR THEME COMMANDS (Robust Single-Slide Overrides)
% -------------------------------------------------------------------------

% 1. Light Mode Command
\newcommand{\LightMode}{
  % 1. Draw White Background (TikZ Overlay)
  \begin{tikzpicture}[remember picture, overlay]
    \fill[white] (current page.south west) rectangle (current page.north east);
  \end{tikzpicture}
  % 2. Force Text Colors to Black
  \setbeamercolor{normal text}{fg=black}
  \setbeamercolor{frametitle}{fg=black}
  \setbeamercolor{itemize item}{fg=black}
  \setbeamercolor{structure}{fg=black}
  \usebeamercolor[fg]{normal text}
}

% 2. Black Mode Command
\newcommand{\BlackMode}{
  \begin{tikzpicture}[remember picture, overlay]
    \fill[black] (current page.south west) rectangle (current page.north east);
  \end{tikzpicture}
  \setbeamercolor{normal text}{fg=white}
  \setbeamercolor{frametitle}{fg=white}
  \setbeamercolor{itemize item}{fg=white}
  \setbeamercolor{structure}{fg=white}
  \usebeamercolor[fg]{normal text}
}

% 3. Brand Mode Command
\newcommand{\BrandMode}{
  \begin{tikzpicture}[remember picture, overlay]
    \definecolor{tempBrand}{HTML}{7D181E}
    \fill[tempBrand] (current page.south west) rectangle (current page.north east);
  \end{tikzpicture}
  \setbeamercolor{normal text}{fg=white}
  \setbeamercolor{frametitle}{fg=white}
  \setbeamercolor{itemize item}{fg=white}
  \setbeamercolor{structure}{fg=white}
  \usebeamercolor[fg]{normal text}
}

% -------------------------------------------------------------------------
% 4. UC BERKELEY THEME (Modern Matte) - FIXED SCOPE
% -------------------------------------------------------------------------

% 1. DEFINE COLORS GLOBALLY (Required for assumptionbox and modes)
\definecolor{BerkeleyBlue}{HTML}{003262}
\definecolor{CaliforniaGold}{HTML}{FDB515}
\definecolor{BerkeleyMatte}{HTML}{0D2D52} 
\definecolor{SoftWhite}{HTML}{F8F9FA}

% 4. UC Berkeley Theme (Modern Matte)
\newcommand{\BerkeleyMode}{
  % 2. Draw Background AND Title
  \begin{tikzpicture}[remember picture, overlay]
    % Background Fill (Covers old title)
    \fill[BerkeleyMatte] (current page.south west) rectangle (current page.north east);
    % Gold Accent Strip
    \fill[CaliforniaGold] (current page.south west) rectangle ($(current page.north west)+(0.15cm,0)$);
    % REDRAW TITLE ON TOP
    \node[anchor=west, white, font=\Large\bfseries] at ($(current page.north west)+(1cm,-1cm)$) {\insertframetitle};
  \end{tikzpicture}
  
  % 3. Set Text Colors
  \setbeamercolor{normal text}{fg=SoftWhite}
  \setbeamercolor{frametitle}{fg=BerkeleyMatte} % Hide original title
  \setbeamercolor{itemize item}{fg=CaliforniaGold} 
  \setbeamercolor{structure}{fg=CaliforniaGold}
  
  \usebeamercolor[fg]{normal text}
}

% 5. UC Berkeley Theme (Academic Light)
\newcommand{\BerkeleyLightMode}{
  % (Color definitions removed here as they are now global)
  
  % 2. Draw Background AND Title
  \begin{tikzpicture}[remember picture, overlay]
    % Background Fill
    \fill[white] (current page.south west) rectangle (current page.north east);
    % Blue Header Bar
    \fill[BerkeleyBlue] (current page.north west) rectangle ($(current page.north east)+(0,-1.4cm)$);
    % Gold Line
    \fill[CaliforniaGold] ($(current page.north west)+(0,-1.4cm)$) rectangle ($(current page.north east)+(0,-1.45cm)$);
    % REDRAW TITLE ON TOP
    \node[anchor=west, white, font=\Large\bfseries] at ($(current page.north west)+(1cm,-0.7cm)$) {\insertframetitle};
  \end{tikzpicture}
  
  % 3. Set Text Colors
  \setbeamercolor{normal text}{fg=black}
  \setbeamercolor{frametitle}{fg=white} 
  \setbeamercolor{itemize item}{fg=BerkeleyBlue}
  \setbeamercolor{structure}{fg=BerkeleyBlue}
  
  \usebeamercolor[fg]{normal text}
}

% -------------------------------------------------------------------------
% NEW: ACADEMIC POWER LAYOUTS
% -------------------------------------------------------------------------

% 1. Statement Layout (Big Impact Transitions)
% Usage: ::: {.layout-statement} Content :::
\newenvironment{layout-statement}{
  \vspace*{\fill}
  \centering
  \Huge \bfseries
  \setbeamercolor{normal text}{fg=CaliforniaGold} % Use Gold for impact
  \usebeamercolor[fg]{normal text}
}{
  \vspace*{\fill}
}

% 2. Dense Data Layout (Regression Tables)
% Usage: ::: {.layout-dense} Table :::
\newenvironment{layout-dense}{
  \vspace*{\fill}
  \tiny                  % Force tiny font
  \setlength\tabcolsep{2pt} % Tighten table columns
  \centering
}{
  \vspace*{\fill}
}

% 3. Split Slide Macro (Vertical Centering)
% Usage: \SplitSlide{Left Content}{Right Content}
\newcommand{\SplitSlide}[2]{
  \begin{columns}[c] % 'c' option forces vertical centering
    \begin{column}{0.48\textwidth}
      #1
    \end{column}
    \begin{column}{0.48\textwidth}
      \centering
      #2
    \end{column}
  \end{columns}
}

% -------------------------------------------------------------------------
% 3. DYNAMIC BACKGROUND & CONTRAST LOGIC
% -------------------------------------------------------------------------
  \definecolor{mainbg}{HTML}{0D2D52}

\setbeamercolor{background canvas}{bg=mainbg}

% Logic to calculate luminance and set text color automatically
\makeatletter
\newcommand{\AutoSetContrast}{
  \extractcolorspec{mainbg}{\mainbgspec}
  \expandafter\convertcolorspec\mainbgspec{rgb}\mainbgrgb
  \def\calculatebrightness##1,##2,##3;{
    \pgfmathparse{0.2126*##1 + 0.7152*##2 + 0.0722*##3}
  }
  \expandafter\calculatebrightness\mainbgrgb;
  \pgfmathparse{\pgfmathresult > 0.5 ? 1 : 0}
  \ifnum\pgfmathresult=1
    \definecolor{maintext}{HTML}{222222}
    \definecolor{dimtext}{HTML}{555555}
  \else
    \definecolor{maintext}{HTML}{FFFFFF}
    \definecolor{dimtext}{HTML}{CCCCCC}
  \fi
}
\makeatother

\AutoSetContrast

\setbeamercolor{normal text}{fg=maintext}
\setbeamercolor{frametitle}{fg=maintext}
\setbeamercolor{title}{fg=maintext}
\setbeamercolor{structure}{fg=maintext} 

% -------------------------------------------------------------------------
% 4. FOOTER CUSTOMIZATION
% -------------------------------------------------------------------------
\setbeamercolor{footline}{fg=dimtext}
\setbeamerfont{footline}{size=\tiny}

\newtoggle{hidesection}
\togglefalse{hidesection}

\setbeamertemplate{footline}{
  \begin{beamercolorbox}[wd=\paperwidth, ht=2.5ex, dp=1.125ex, leftskip=.3cm, rightskip=.3cm]{footline}
    \iftoggle{hidesection}{}{
      \insertsectionhead
    }
    \hfill
    \insertframenumber
  \end{beamercolorbox}
}
\setbeamertemplate{navigation symbols}{}

% -------------------------------------------------------------------------
% 5. CUSTOM COMMANDS & ENVIRONMENTS (FIXED FOR STABILITY)
% -------------------------------------------------------------------------
\NewDocumentCommand{\navbutton}{ O{} O{fill=gray} m }{
  \hyperlink{#1}{
    \tikz[baseline=(node.base)]{
      \node(node)[
        anchor=base,
        rounded corners=3pt,
        inner sep=5pt,
        text=white,
        font=\small\bfseries,
        #2
      ]{#3};
    }
  }
}

\newenvironment{layout-math}{
  \vspace*{\fill}
  \centering
  \begin{minipage}{\linewidth}
  \centering
}{
  \end{minipage}
  \vspace*{\fill}
}

% SAFE Layout for Tables (Removed risky \makebox)
\newenvironment{layout-table}{
  \vspace*{\fill}
  \begin{center}
}{
  \end{center}
  \vspace*{\fill}
}

% SAFE Layout for Full Figures (Uses LRBOX to capture content first)
\newsavebox{\fullfigbox}
\newenvironment{layout-full-figure}{
  \begin{lrbox}{\fullfigbox}
    \begin{minipage}{\paperwidth}
      \centering
}{
    \end{minipage}
  \end{lrbox}
  \begin{tikzpicture}[remember picture, overlay]
    \node[anchor=center, inner sep=0pt] at (current page.center) {
      \usebox{\fullfigbox}
    };
  \end{tikzpicture}
}

\newcommand{\hidefootersection}{\global\toggletrue{hidesection}}
\newcommand{\showfootersection}{\global\togglefalse{hidesection}}


% -------------------------------------------------------------------------
% CUSTOM BERKELEY CALLOUT BOX
% -------------------------------------------------------------------------
\newtcolorbox{assumptionbox}[1]{
  enhanced,
  colback=SoftWhite,       % White-ish background
  colframe=CaliforniaGold, % Gold Border
  coltitle=SoftWhite,      % White Title Text
  colbacktitle=BerkeleyMatte, % Berkeley Blue Title Background
  title={#1},
  fonttitle=\bfseries\small,
  arc=2pt,                 % Slight rounded corners
  boxrule=1pt,             % Thin elegant border
  left=5pt, right=5pt, top=5pt, bottom=5pt, % Padding
  titlerule=0pt,
  width=\linewidth
}

% -------------------------------------------------------------------------
% 6. DOCUMENT BODY
% -------------------------------------------------------------------------
\title{Universal DiD Analysis: Texas UST Insurance Reform}
\author{Kaleb K.}

\begin{document}

\frame{\titlepage}

\begin{frame}{Goals for Today's Talk}
\phantomsection\label{goals-for-todays-talk}
\BerkeleyLightMode

\textbf{1. Toy Model: Is the Focus Right?}

\begin{itemize}
\tightlist
\item
  Does the simple framework clearly capture the behavior I can measure
  (tank closures)?
\item
  Thoughts on how to sharpen the premium \(\rightarrow\) closure
  mechanism?
\end{itemize}

\textbf{2. Structural Model: What's Missing?}

\begin{itemize}
\tightlist
\item
  Is the DDC setup complete, or are key state variables/parameters
  omitted?
\item
  \textbf{Specific question:} How do I incorporate Marcus (2021) health
  externalities without estimating damages myself? Just use their
  estimates as inputs?
\item
  Best ways to push this forward?
\end{itemize}
\end{frame}

\begin{frame}{Event Study: Empirical Specification}
\phantomsection\label{event-study-empirical-specification}
\BerkeleyLightMode

\begin{assumptionbox}{Dynamic Treatment Effects}
$$Y_{it} = \sum_{\tau \neq \tau_{ref}} \beta_\tau (\text{Texas}_i \times \mathbb{I}(t = \tau)) + \gamma X_{it} + \alpha_i + \delta_{ct} + \varepsilon_{it}$$
\end{assumptionbox}

\begin{itemize}
\tightlist
\item
  \(Y_{it}\): Any Tank Closure (0/1)
\item
  \(\text{Texas}_i \times \mathbb{I}(t = \tau)\): Event-time indicators,
  1999 reform
\item
  \(X_{it}\): Mandate controls (Release Detection, Spill/Overfill)
\item
  \(\alpha_i\): Facility FE; \(\quad\) \(\delta_{ct}\): Cell \(\times\)
  Year FE \(\;\) (Cell = wall \(\times\) fuel \(\times\) age bin)
\end{itemize}

\textbf{Identification:} \(\delta_{ct}\) relaxes parallel trends to
\textbf{within-cell}, restricting the counterfactual to facilities with
identical technological vintage and risk profile.
\end{frame}

\begin{frame}{}
\phantomsection\label{section}
\BerkeleyLightMode

\hidefootersection

\begin{layout-full-figure}

\includegraphics[width=1\textwidth,height=\textheight]{//ucbare2.are.berkeley.edu/kalebkja/ust_ins_move_to_github/Output/Figures/Figure_ES_Sec1_Ref-1_Naive_AnyClosure_MainPaper.pdf}

\end{layout-full-figure}

\showfootersection
\end{frame}

\begin{frame}{}
\phantomsection\label{section-1}
\BerkeleyLightMode

\hidefootersection

\begin{layout-full-figure}

\includegraphics[width=1\textwidth,height=\textheight]{//ucbare2.are.berkeley.edu/kalebkja/ust_ins_move_to_github/Output/Figures/Figure_ES_Sec2_Ref-2_Cell_AnyClosure_MainPaper.pdf}

\end{layout-full-figure}

\showfootersection
\end{frame}

\begin{frame}{}
\phantomsection\label{section-2}
\BerkeleyLightMode

\hidefootersection

\begin{layout-full-figure}

\includegraphics[width=1\textwidth,height=\textheight]{//ucbare2.are.berkeley.edu/kalebkja/ust_ins_move_to_github/Output/Figures/Figure_ES_Sec3_Ref-1_Cell_AnyClosure_MainPaper.pdf}

\end{layout-full-figure}

\showfootersection
\end{frame}

\begin{frame}{Static DiD: Extensive Margin (Any Closure)}
\phantomsection\label{static-did-extensive-margin-any-closure}
\BerkeleyLightMode

\begin{assumptionbox}{Baseline Static Specification}
$$Y_{it} = \beta^{DiD} (\text{Texas}_i \times \text{Post}_t) + \gamma X_{it} + \alpha_i + \delta_{ct} + \varepsilon_{it}$$
\end{assumptionbox}

\begin{itemize}
\tightlist
\item
  \(Y_{it}\): Any Tank Closure (0/1)
\item
  \(\text{Texas}_i \times \text{Post}_t\): Binary treatment
  (\(t \geq 1999\))
\item
  \(\alpha_i\), \(\delta_{ct}\), \(X_{it}\): Facility FE, Cell
  \(\times\) Year FE, mandate controls
\end{itemize}

\textbf{Identification:} \(\beta^{DiD}\) estimates the average closure
probability shift; \(\delta_{ct}\) absorbs endogenous retirement waves
from RCRA deadlines hitting specific tank vintages.
\end{frame}

\begin{frame}{}
\phantomsection\label{section-3}
\BerkeleyLightMode

\begin{layout-dense}
\begin{table}[htbp]
\centering
\caption{Table 1: Do Facilities Close More Tanks? Effect of the Texas Insurance Reform on Any Tank Closure Activity.}
\label{tab:t1_any_closure}
\begin{tabular}{lcccccc}
\toprule
 & (1) & (2) & (3) & (4) & (5) & (6) \\
 & Naive & +Mandates & +Cell$\times$Yr & CEM+Mandates & CEM+Cell$\times$Yr & CEM+2D-Cell$\times$Yr \\
\midrule
TX $\times$ Post & 0.0163*** & 0.0163*** & 0.0147*** & 0.0159*** & 0.0130** & 0.0158*** \\
                  & (0.0048) & (0.0048) & (0.0049) & (0.0035) & (0.0047) & (0.0042) \\
$N$               & 3,559,551 & 3,559,551 & 2,614,051 & 2,605,863 & 2,605,834 & 2,605,863 \\
Control mean      & \multicolumn{6}{c}{0.0326} \\
\midrule
Facility FE              & Yes & Yes & Yes & Yes & Yes & Yes \\
Year FE                  & Yes & Yes & Yes & Yes & Yes & Yes \\
Cell$\times$Year FE     & No  & No  & 3D  & No  & 3D  & 2D  \\
CEM weights              & No  & No  & No  & Yes & Yes & Yes \\
Mandate controls         & No  & Yes & Yes & Yes & Yes & Yes \\
Sample & \multicolumn{3}{c}{Full incumbents (no matching)} & \multicolumn{3}{c}{Reform-CEM matched} \\
\bottomrule
\end{tabular}
\begin{minipage}{\linewidth}
\small
\textit{Notes:} LPM coefficients. Dependent variable: Any Tank Closure (0/1). Cols (1)--(3) use the full sample of incumbent facilities with no matching or reweighting. Cols (4)--(6) use reform-survivor CEM matched facilities (alive Dec 22 1998, matched on tank count $\times$ SW count $\times$ fuel type $\times$ oldest tank age bin). 3D Cell$\times$Year FE: make\_model\_fac (wall$\times$fuel$\times$age) $\times$ year. 2D Cell$\times$Year FE: fac\_wall$\times$fac\_fuel$\times$year (drops age bin; robustness). Col (5) is the primary specification used in all subsequent tables. Cluster-robust SEs by state (17 clusters). $^{*}p<0.10$, $^{**}p<0.05$, $^{***}p<0.01$.
\end{minipage}
\end{table}

\end{layout-dense}
\end{frame}

\begin{frame}{Closure Composition: What Do Firms Do When They Close?}
\phantomsection\label{closure-composition-what-do-firms-do-when-they-close}
\BerkeleyLightMode

\begin{assumptionbox}{Descriptive Conditional Specification}
$$\mathbb{E}[Y_{it}^k \mid \text{Closure}_{it} = 1] = \beta_k (\text{Texas}_i \times \text{Post}_t) + \gamma X_{it} + \alpha_i + \delta_{ct}$$
\end{assumptionbox}

\begin{itemize}
\tightlist
\item
  \(Y_{it}^k \in \{\text{Exit}, \text{Upgrade}, \text{Shrink}\}\):
  Closure type, conditional on closing
\item
  \(\beta_k\): \textbf{Descriptive only} --- conditions on the
  endogenous decision to close
\end{itemize}

\textbf{Note:} \(\delta_{ct}\) ensures compositional shifts aren't
driven by aging single-wall cohorts retiring on schedule.
\end{frame}

\begin{frame}{}
\phantomsection\label{section-4}
\BerkeleyLightMode

\begin{layout-dense}
\begin{table}[htbp]
\centering
\caption{Table 2b: What Type of Closure? Descriptive Composition of Tank Closure Events Among Facilities That Closed at Least One Tank.}
\label{tab:t2b_closure_conditional}
\begin{tabular}{lcccc}
\toprule
Closure Type & TX $\times$ Post & (SE) & Bootstrap $p$ & Control Mean \\
\midrule
Facility Closes Entirely & 0.0718*** & (0.0200) &   NA & 0.0065 \\
Facility Upgrades a Tank & 0.0637*** & (0.0129) &   NA & 0.4549 \\
Facility Shrinks Tank Stock & -0.0643*** & (0.0122) &   NA & 0.5562 \\
\midrule
$N$ (closing fac-years) & \multicolumn{4}{c}{33,966} \\
Cell$\times$Year FE & \multicolumn{4}{c}{3D (make\_model\_fac $\times$ year)} \\
CEM weights          & \multicolumn{4}{c}{Yes} \\
Mandate controls     & \multicolumn{4}{c}{Yes} \\
Sample               & \multicolumn{4}{c}{Reform-CEM, any\_closure $= 1$} \\
\bottomrule
\end{tabular}
\begin{minipage}{\linewidth}
\small
\textit{Notes:} LPM coefficients. Sample restricted to reform-CEM facility-years in which at least one tank closure occurred (any\_closure $= 1$). Each row describes the composition of closure events: conditional on closing a tank that year, was the event a full facility exit, a technology upgrade, or a tank stock reduction? \textit{These estimates are descriptive -- they condition on the endogenous decision to close a tank and should not be interpreted as causal effects of the reform on closure type.} For unconditional causal estimates see Table 2. Bootstrap p-values are from a wild cluster score bootstrap (Kline-Santos 2012) with Rademacher weights, B = 9999, clustering by state. Cluster-robust SEs by state. $^{*}p<0.10$, $^{**}p<0.05$, $^{***}p<0.01$.
\end{minipage}
\end{table}

\end{layout-dense}
\end{frame}

\begin{frame}{Portfolio Capital Dynamics: How Does the Stock Change?}
\phantomsection\label{portfolio-capital-dynamics-how-does-the-stock-change}
\BerkeleyLightMode

\begin{assumptionbox}{Descriptive Portfolio Outcomes (unconditional)}
$$Y_{it}^{port} = \beta^{port} (\text{Texas}_i \times \text{Post}_t) + \gamma X_{it} + \alpha_i + \delta_{ct} + \varepsilon_{it}$$
\end{assumptionbox}

\begin{itemize}
\tightlist
\item
  \(Y_{it}^{port}\): Net tank change, total capacity change,
  \(\text{pct\_sw}\) (single-wall share)
\item
  \textbf{Unconditional} --- unlike Table 2b, does not condition on
  closing
\end{itemize}

\textbf{Note:} \(\beta^{port}\) describes the net shift in aggregate
capital stock; \(\delta_{ct}\) absorbs cohort-driven retirement trends
within tank vintage \(\times\) year cells.
\end{frame}

\begin{frame}{}
\phantomsection\label{section-5}
\BerkeleyLightMode

\begin{layout-dense}
\begin{table}[htbp]
\centering
\caption{Table 3: What Do Replacements Do to the Portfolio? Effect of the Texas Insurance Reform on Facility Portfolio Composition.}
\label{tab:t3_portfolio}
\resizebox{\textwidth}{!}{
\begin{tabular}{lcccc}
\toprule
Outcome & (1) Year FE & (2) Cell$\times$Yr & (3) 2D-Cell$\times$Yr & Ctrl Mean \\
\midrule
SW Tank Removed, DW Tank Installed (0/1) & 0.0018*** & 0.0012* & 0.0016** & 0.0030 \\
 & (0.0005) & (0.0006) & (0.0007) &  \\[0.3em]
Any DW Tank Installed This Year (0/1) & 0.0006 & -0.0009 & -0.0004 & 0.0138 \\
 & (0.0026) & (0.0010) & (0.0012) &  \\[0.3em]
Net Change in Tank Count (continuous) & -0.0314*** & -0.0260** & -0.0304*** & -0.0263 \\
 & (0.0094) & (0.0119) & (0.0101) &  \\[0.3em]
Change in Total Capacity, gallons (continuous) & -316.0850*** & -288.3101*** & -308.6747*** & -98.2379 \\
 & (60.7569) & (67.9809) & (61.9043) &  \\[0.3em]
Single-Wall Tanks at End of Year (continuous) & -0.0502 & -0.0675 & -0.0550 & 2.6410 \\
 & (0.0548) & (0.0478) & (0.0494) &  \\[0.3em]
Double-Wall Tanks at End of Year (continuous) & -0.0364 & -0.0203 & -0.0301** & 0.1955 \\
 & (0.0250) & (0.0137) & (0.0138) &  \\[0.3em]
Share of Portfolio That Is Single-Wall (0--1) & -0.0047 & -0.0110 & -0.0078 & 0.9082 \\
 & (0.0108) & (0.0076) & (0.0076) &  \\[0.3em]
\midrule
$N$ & \multicolumn{4}{c}{2,605,834} \\
Facility FE            & Yes & Yes & Yes & \\
Year FE                & Yes & Yes & Yes & \\
Cell$\times$Year FE   & No  & 3D  & 2D  & \\
CEM weights            & Yes & Yes & Yes & \\
Mandate controls       & Yes & Yes & Yes & \\
Sample & \multicolumn{4}{c}{Reform-CEM matched, all facility-years} \\
\bottomrule
\end{tabular}
}
\begin{minipage}{\linewidth}
\small
\textit{Notes:} OLS coefficients. Unconditional estimates on all reform-CEM facility-years. Binary outcomes (rows 1--2) estimated by LPM; continuous outcomes (rows 3--7) by OLS. The reform shifts portfolio composition through the replacement channel documented in Table 2 Panel B, but these regressions do not condition on replacement occurring and are causally interpretable. 3D Cell$\times$Year FE: make\_model\_fac (wall$\times$fuel$\times$age) $\times$ year. 2D Cell$\times$Year FE: fac\_wall$\times$fac\_fuel$\times$year (robustness: drops age bin from cell definition). Cluster-robust SEs by state (17 clusters). $^{*}p<0.10$, $^{**}p<0.05$, $^{***}p<0.01$.
\end{minipage}
\end{table}

\end{layout-dense}
\end{frame}

\begin{frame}{DiD: Leak Discovery Incidence (LUST)}
\phantomsection\label{did-leak-discovery-incidence-lust}
\BerkeleyLightMode

\begin{assumptionbox}{Total and Decomposed LUST Incidence}
$$\text{Leak}_{it}^m = \beta^m (\text{Texas}_i \times \text{Post}_t) + \gamma X_{it} + \alpha_i + \delta_{ct} + \varepsilon_{it}$$
\end{assumptionbox}

\begin{itemize}
\tightlist
\item
  \(\text{Leak}_{it}^{total}\): Any leak discovery in year \(t\)
\item
  \(\text{Leak}_{it}^{back}\): Background leak (routine monitoring)
\item
  \(\text{Leak}_{it}^{inspect}\): Inspection-triggered (within \(X\)
  days of closure)
\end{itemize}

\textbf{Identification:} \(\delta_{ct}\) absorbs age-driven failure
rates; \(\beta^{back}\) vs.~\(\beta^{inspect}\) isolates whether the
reform accelerates \emph{discovery} of latent contamination versus
altering the true failure rate.
\end{frame}

\begin{frame}{}
\phantomsection\label{section-6}
\BerkeleyLightMode

\begin{layout-dense}
\begin{table}[htbp]
\centering
\caption{LUST Table A: Does the Reform Change the Rate of Leak Discoveries? Effect of the Texas Insurance Reform on Total LUST Incidence.}
\label{tab:lust_a_total}
\begin{tabular}{lccc}
\toprule
 & (1) Year FE & (2) Cell$\times$Yr & (3) 2D-Cell$\times$Yr \\
\midrule
TX $\times$ Post & -0.0107*** & -0.0113*** & -0.0108*** \\
                  & (0.0030) & (0.0033) & (0.0031) \\
$N$               & 2,605,863 & 2,605,834 & 2,605,863 \\
Control mean      & \multicolumn{3}{c}{0.0193} \\
\midrule
Facility FE            & Yes & Yes & Yes \\
Year FE                & Yes & Yes & Yes \\
Cell$\times$Year FE   & No  & 3D  & 2D  \\
CEM weights            & Yes & Yes & Yes \\
Mandate controls       & Yes & Yes & Yes \\
Sample & \multicolumn{3}{c}{Reform-CEM matched} \\
\bottomrule
\end{tabular}
\begin{minipage}{\linewidth}
\small
\textit{Notes:} LPM coefficients. Dependent variable: Any Leak Discovery This Year (0/1). Total leak discovery includes leaks found through any channel (routine monitoring, complaints, and closure inspections). See LUST Table B for decomposition by discovery channel. Cluster-robust SEs by state (17 clusters). $^{*}p<0.10$, $^{**}p<0.05$, $^{***}p<0.01$.
\end{minipage}
\end{table}

\end{layout-dense}
\end{frame}

\begin{frame}{}
\phantomsection\label{section-7}
\BerkeleyLightMode

\begin{layout-dense}
\begin{table}[htbp]
\centering
\caption{LUST Table B: Background vs.\ Inspection-Triggered Leak Discoveries. Decomposition of LUST Incidence by Discovery Channel and Inspection Window Definition.}
\label{tab:lust_b_decomp}
\begin{tabular}{lcccc}
\toprule
Window & (1) Year FE & (2) Cell$\times$Yr & (3) 2D-Cell$\times$Yr & Ctrl Mean \\
\midrule
\multicolumn{5}{l}{\textit{Panel A: Background Leaks}} \\
\multicolumn{5}{l}{\small\textit{Leak reported with no tank closure within the window. Reflects routine monitoring or complaint-driven discovery. }} \\
\multicolumn{5}{l}{\small\textit{An increase here indicates more genuine operational failures or increased non-closure detection.}} \\[0.3em]
Primary (0--60 day window) & -0.0065** & -0.0067** & -0.0065** & 0.0157 \\
 & (0.0024) & (0.0025) & (0.0025) &  \\[0.2em]
Narrow (0--30 day window) & -0.0070** & -0.0072** & -0.0070** & 0.0163 \\
 & (0.0025) & (0.0026) & (0.0025) &  \\[0.2em]
Regulatory (0--45 day window) & -0.0067** & -0.0069** & -0.0067** & 0.0159 \\
 & (0.0025) & (0.0025) & (0.0025) &  \\[0.2em]
Wide (0--90 day window) & -0.0062** & -0.0063** & -0.0062** & 0.0154 \\
 & (0.0024) & (0.0025) & (0.0024) &  \\[0.2em]
\midrule
\multicolumn{5}{l}{\textit{Panel B: Inspection-Triggered Leaks}} \\
\multicolumn{5}{l}{\small\textit{Leak reported within the window around a tank closure. Closures require regulatory inspection; inspections reveal }} \\
\multicolumn{5}{l}{\small\textit{latent contamination that already existed. An increase here reflects faster discovery of pre-existing leaks -- an indirect environmental benefit.}} \\[0.3em]
Primary (0--60 day window) & -0.0043** & -0.0048** & -0.0044** & 0.0037 \\
 & (0.0016) & (0.0017) & (0.0016) &  \\[0.2em]
Narrow (0--30 day window) & -0.0038** & -0.0041** & -0.0038** & 0.0030 \\
 & (0.0014) & (0.0015) & (0.0014) &  \\[0.2em]
Regulatory (0--45 day window) & -0.0041** & -0.0045** & -0.0042** & 0.0034 \\
 & (0.0015) & (0.0016) & (0.0015) &  \\[0.2em]
Wide (0--90 day window) & -0.0047** & -0.0051** & -0.0047** & 0.0041 \\
 & (0.0018) & (0.0019) & (0.0018) &  \\[0.2em]
\midrule
$N$ & \multicolumn{4}{c}{2,605,834} \\
Facility FE            & Yes & Yes & Yes & \\
Year FE                & Yes & Yes & Yes & \\
Cell$\times$Year FE   & No  & 3D  & 2D  & \\
CEM weights            & Yes & Yes & Yes & \\
Mandate controls       & Yes & Yes & Yes & \\
Sample & \multicolumn{4}{c}{Reform-CEM matched} \\
\bottomrule
\end{tabular}
\begin{minipage}{\linewidth}
\small
\textit{Notes:} LPM coefficients. Each row reports estimates under a different window definition for classifying a leak report as inspection-triggered or background. The window is the number of days between a tank closure and a leak report at the same facility. Panel A leaks have no closure within the window; Panel B leaks have a closure within the window. The 60-day primary window is the main specification; 30-day, 45-day, and 90-day windows test sensitivity to the threshold. 3D Cell$\times$Year FE: make\_model\_fac $\times$ year. 2D Cell$\times$Year FE: fac\_wall$\times$fac\_fuel$\times$year. Cluster-robust SEs by state (17 clusters). $^{*}p<0.10$, $^{**}p<0.05$, $^{***}p<0.01$.
\end{minipage}
\end{table}

\end{layout-dense}
\end{frame}

\begin{frame}{}
\phantomsection\label{section-8}
\BerkeleyMode

Dynamic Optimal Stopping Model
\end{frame}

\begin{frame}{Model Setup}
\phantomsection\label{model-setup}
\BerkeleyLightMode

\begin{assumptionbox}{Value Function}
$$V(s_t) = \max_{d_t \in \{0, 1, 2\}} \left[ \pi(s_t, d_t) + \beta \mathbb{E}[V(s_{t+1}) \mid s_t, d_t] \right]$$
\end{assumptionbox}

\textbf{Choice Set (\(d_t\)):}

\begin{itemize}
\tightlist
\item
  \(d_t = 0\): \textbf{Continue} --- status quo utility, stochastic age
  transition
\item
  \(d_t = 1\): \textbf{Retrofit} --- incur fixed cost \(K_R\), reset
  regulatory state
\item
  \(d_t = 2\): \textbf{Exit} --- claim terminal scrap value \(K_E\),
  permanent absorption
\end{itemize}
\end{frame}

\begin{frame}{4-Parameter Baseline: State Space \& Bellman Equation}
\phantomsection\label{parameter-baseline-state-space-bellman-equation}
\BerkeleyLightMode

\begin{assumptionbox}{Dynamic Programming Framework — Homogeneous Costs}
$$V(s_t) = \max_{d_t \in \{0, 1, 2\}} \left[ u(s_t, d_t) + \beta \int V(s_{t+1}) dP(s_{t+1} \mid s_t, d_t) \right]$$
\end{assumptionbox}

\textbf{State Space:} \(s_t = (a_t, \epsilon_t)\) --- tank age and
i.i.d. Type I EV profitability shock.

Identical compliance cost structures assumed across all facilities.
\end{frame}

\begin{frame}{4-Parameter Baseline: Utility Specification}
\phantomsection\label{parameter-baseline-utility-specification}
\BerkeleyLightMode

\(\theta = (\alpha, \gamma, K_R, K_E)\)

\begin{itemize}
\tightlist
\item
  \textbf{Continue (\(d_t = 0\)):}
  \(u = \alpha - \gamma a_t + \epsilon_{0t}\) \(\quad\) \emph{(base
  profit, age penalty)}
\item
  \textbf{Retrofit (\(d_t = 1\)):} \(u = \alpha - K_R + \epsilon_{1t}\)
  \(\quad\) \emph{(fixed upgrade cost, resets \(a_{t+1}=0\))}
\item
  \textbf{Exit (\(d_t = 2\)):} \(u = K_E + \epsilon_{2t}\) \(\quad\)
  \emph{(terminal scrap value, absorbing)}
\end{itemize}
\end{frame}

\begin{frame}{4-Parameter Baseline: Transition \& Likelihood}
\phantomsection\label{parameter-baseline-transition-likelihood}
\BerkeleyLightMode

\textbf{Transitions:} Continue \(\Rightarrow a_t+1\); \(\;\) Retrofit
\(\Rightarrow a=0\); \(\;\) Exit \(\Rightarrow\) absorbing.

\[P(d_t \mid a_t ; \theta) = \frac{\exp(v(a_t, d_t))}{\sum_{j=0}^{2} \exp(v(a_t, j))}\]

Conditional independence of \(\epsilon_t\) yields closed-form logistic
choice probabilities. MLE over observed panel of facility actions.
\end{frame}

\begin{frame}{Identification: 4-Parameter Baseline}
\phantomsection\label{identification-4-parameter-baseline}
\BerkeleyLightMode

Baseline Identification Strategy

\begin{itemize}
\tightlist
\item
  \textbf{Age Gradient:} Cross-sectional variation in tank age at the
  time of the policy shock.
\item
  \textbf{Deadline Proximity:} Sharp discontinuity in exit/retrofit
  hazard rates preceding the 1998 EPA deadline.
\item
  \textbf{Assumption:} Homogeneous compliance costs across all incumbent
  facilities.
\end{itemize}
\end{frame}

\begin{frame}{4-Parameter Model Estimates}
\phantomsection\label{parameter-model-estimates}
\BerkeleyLightMode

\begin{layout-dense}
% Auto-generated by 04f_Replacement_Refit_Extended2000_AM_SE.R
% AM (2002) corrected SEs in (parentheses); bootstrap SEs deferred.
\begin{tabular}{lcc}
\hline
 & Observed & Extended (2000+) \\
 & TX 2006+, controls 1999+ & TX 2000+, controls 2000+ \\
\hline
$\kappa_{\mathrm{exit}}$ (\$) & \$261,222 & \$257,702 \\
        & (\$306) & (\$310) \\
$K$ (\$) & \$51,779 & \$52,658 \\
        & (\$142) & (\$146) \\
$\gamma_{\mathrm{price}}$ & -0.293 & -1.142 \\
        & (0.049) & (0.048) \\
$\gamma_{\mathrm{risk}}$ & 0.164 & 0.153 \\
        & (0.003) & (0.003) \\
\hline
$\log L$ & -268,481 & -247,614 \\
$N$ obs  & 2,282,735 & 2,300,780 \\
\hline
\multicolumn{3}{l}{\footnotesize SEs are AM (2002) profile-likelihood;
  bootstrap SEs deferred for a future pass.}
\end{tabular}

\end{layout-dense}
\end{frame}

\begin{frame}{8-Parameter Heterogeneity: Expanded State Space}
\phantomsection\label{parameter-heterogeneity-expanded-state-space}
\BerkeleyLightMode

\begin{assumptionbox}{Relaxes uniform compliance costs across wall type and fuel regime}
$$V(a_t, W_i, R_i) = \max_{d_t} \left[ u(a_t, W_i, R_i, d_t) + \beta \mathbb{E}[V(a_{t+1}, W_i, R_i) \mid a_t, d_t] \right]$$
\end{assumptionbox}

\textbf{Expanded State:} \(s_t = (a_t, W_i, R_i, \epsilon_t)\)

\begin{itemize}
\tightlist
\item
  \(W_i \in \{\text{SW}, \text{DW}\}\): Tank wall type (single-wall
  vs.~double-wall)
\item
  \(R_i \in \{\text{FF}, \text{RB}\}\): Fuel regime (fossil fuel
  vs.~reformulated blend)
\end{itemize}
\end{frame}

\begin{frame}{8-Parameter Heterogeneity: Utility Specification}
\phantomsection\label{parameter-heterogeneity-utility-specification}
\BerkeleyLightMode

\(\theta_{het} = (\kappa_{\mathrm{SW}}, \kappa_{\mathrm{DW}}, K_{\mathrm{SW}}, K_{\mathrm{DW}}, \gamma_{\mathrm{price,FF}}, \gamma_{\mathrm{price,RB}}, \gamma_{\mathrm{risk,FF}}, \gamma_{\mathrm{risk,RB}})\)

\begin{itemize}
\tightlist
\item
  \textbf{Continue (\(d_t = 0\)):}
  \(u = -\gamma_{\mathrm{price},R_i} \cdot p_t - \gamma_{\mathrm{risk},R_i} \cdot a_t + \epsilon_{0t}\)
  \(\quad\) \emph{(price and risk sensitivity by fuel regime)}
\item
  \textbf{Retrofit (\(d_t = 1\)):} \(u = -K_{W_i} + \epsilon_{1t}\)
  \(\quad\) \emph{(upgrade cost by wall type)}
\item
  \textbf{Exit (\(d_t = 2\)):} \(u = \kappa_{W_i} + \epsilon_{2t}\)
  \(\quad\) \emph{(scrap value by wall type)}
\end{itemize}
\end{frame}

\begin{frame}{Identification: 8-Parameter Heterogeneity}
\phantomsection\label{identification-8-parameter-heterogeneity}
\BerkeleyLightMode

Heterogeneous Identification Strategy

\begin{itemize}
\tightlist
\item
  \textbf{\(\kappa_{\mathrm{SW}}\) vs.~\(\kappa_{\mathrm{DW}}\):} Exit
  scrap values differ by wall type --- single-wall faces higher
  remediation liability on exit.
\item
  \textbf{\(K_{\mathrm{SW}}\) vs.~\(K_{\mathrm{DW}}\):} Retrofit costs
  differ by wall type --- upgrading SW infrastructure is more
  capital-intensive.
\item
  \textbf{\(\gamma_{\mathrm{price},FF}\)
  vs.~\(\gamma_{\mathrm{price},RB}\):} Price sensitivity identified by
  differential premium exposure across fuel regimes.
\item
  \textbf{\(\gamma_{\mathrm{risk},FF}\)
  vs.~\(\gamma_{\mathrm{risk},RB}\):} Risk/age penalty identified by
  divergence in age-driven closure hazards across regimes.
\end{itemize}

\textbf{Punchline:} Decomposes whether reform response was driven by
wall-type capital costs or fuel-regime risk pricing.
\end{frame}

\begin{frame}{8-Parameter Model Estimates}
\phantomsection\label{parameter-model-estimates-1}
\BerkeleyLightMode

\begin{layout-dense}
% Auto-generated by 04h_Replacement_8param_Estimation.R
% AM (2002) profile-likelihood SEs in parentheses; 95% Wald CIs in brackets.
% Bootstrap SEs and MC robustness on 8-param spec deferred.
\begin{tabular}{lc}
\hline
 & Observed (headline) \\
 & TX 2006+, controls 1999+ \\
\hline
$\kappa_{\mathrm{SW}}$ & \$236,795 \\
        & (\$716) \\
        & [\$235,392, \$238,198] \\
$\kappa_{\mathrm{DW}}$ & \$211,775 \\
        & (\$856) \\
        & [\$210,098, \$213,452] \\
$K_{\mathrm{SW}}$ & \$58,085 \\
        & (\$212) \\
        & [\$57,670, \$58,501] \\
$K_{\mathrm{DW}}$ & \$43,301 \\
        & (\$244) \\
        & [\$42,823, \$43,779] \\
$\gamma_{\mathrm{price,FF}}$ & -14.382 \\
        & (0.293) \\
        & [-14.957, -13.808] \\
$\gamma_{\mathrm{price,RB}}$ & -7.485 \\
        & (0.151) \\
        & [-7.781, -7.188] \\
$\gamma_{\mathrm{risk,FF}}$ & 0.065 \\
        & (0.005) \\
        & [0.055, 0.075] \\
$\gamma_{\mathrm{risk,RB}}$ & -0.083 \\
        & (0.009) \\
        & [-0.100, -0.066] \\
\hline
$\log L$ & -261,770 \\
$N$ obs  & 2,282,735 \\
\hline
\multicolumn{2}{l}{\footnotesize Point estimates with AM (2002)
  profile-likelihood SEs in parentheses and 95\% CIs in brackets.
  MC robustness battery on 8-param spec deferred.}
\end{tabular}

\end{layout-dense}
\end{frame}

\end{document}
